\chapter{Statistical Modelling and Results}

\section{Modelling Strategy}

The exploratory analysis revealed substantial variation in arrest outcomes
across demographic, temporal, and geographical characteristics. However,
these descriptive comparisons are unconditional: differences observed for one
variable may partly reflect the different composition of the stops with respect
to the remaining predictors.

The modelling stage therefore moves from descriptive comparisons to the
conditional probability

$$
P(\texttt{arrested}=1 \mid X,\text{ recorded vehicle stop}),
$$

where $X$ contains the observed demographic, temporal, and
geographical/organizational characteristics associated with a stop.

Two objectives are pursued simultaneously. The first is
\textbf{interpretative}: we investigate which characteristics remain associated
with arrest when the other observed predictors are considered simultaneously.
The second is \textbf{predictive}: we evaluate how much information these
characteristics contain for distinguishing stops that resulted in an arrest from
those that did not.

The same general modelling strategy is applied independently to the
two-wheeled and four-wheeled vehicle datasets. Four complementary models are
considered:

\begin{itemize}
\item Logistic Regression provides the main interpretable conditional model;
\item LASSO tests whether the main signals survive coefficient shrinkage and
a more parsimonious representation;
\item Generalized Additive Models relax the rigid functional representation
of age and hour while preserving an additive structure;
\item Naive Bayes provides a substantially different, generative view of the
classification problem.
\end{itemize}

The models should therefore not be interpreted as four isolated attempts at
classification. Instead, each model addresses a different aspect of the same
research problem.

% =============================================================================
% LOGISTIC REGRESSION
% =============================================================================

\section{Logistic Regression}

\subsection{Why Logistic Regression?}

Logistic Regression provides the natural starting point for the modelling
analysis because the response \texttt{arrested} is binary and the model directly
estimates the conditional probability of an arrest.

For a vector of predictors $X$,

$$
\log\left(
\frac{P(\texttt{arrested}=1\mid X)}
     {1-P(\texttt{arrested}=1\mid X)}
\right)
=
\beta_0+\beta_1X_1+\cdots+\beta_pX_p.
$$

Its main advantage in this project is interpretability. Each coefficient
describes how the conditional log-odds of arrest differ when one predictor
changes while the remaining observed characteristics are held fixed. For
categorical variables, exponentiating a coefficient produces an odds ratio
relative to the corresponding reference category.

This makes Logistic Regression particularly useful for moving beyond the raw
differences identified during the EDA. The relevant question is no longer only
whether two groups have different observed arrest rates, but whether that
difference remains visible after adjustment for the other available
characteristics.

\subsection{Conditional Associations}

The fitted models confirm that several of the strongest patterns observed
during the EDA remain visible after simultaneous adjustment.

Time of day is one of the clearest examples. In both datasets, the conditional
odds of arrest decline sharply during the morning hours relative to midnight.
For two-wheeled vehicles, the minimum occurs at approximately 07:00, with an
odds ratio of only $0.035$. For four-wheeled vehicles, the corresponding
odds ratio is $0.160$. In both cases, the effects progressively return toward
the midnight reference during the afternoon and evening.

Geography also remains relevant. Using the Bronx as reference, all estimated
borough coefficients are negative in both datasets. For two-wheeled vehicles,
Manhattan North and Manhattan South have odds ratios of $0.332$ and
$0.344$, respectively. In the four-wheeled dataset the corresponding
geographical differences are smaller, although Manhattan South
($OR=0.554$) and Queens North ($OR=0.558$) remain substantially below
the Bronx reference.

The demographic coefficients provide an especially important connection with
the descriptive analysis. With Male as the reference sex category, Female has

$$
OR_{\text{Female}}=0.300
$$

for two-wheeled vehicles and

$$
OR_{\text{Female}}=0.408
$$

for four-wheeled vehicles. Thus, the lower arrest proportions observed for
female motorists in the EDA remain visible after conditioning on the remaining
included predictors.

The contrast across reported ethnicity is even more notable. Using White as
the reference category, the two-wheeled model estimates

$$
OR_{\text{Black}}=1.451,
\qquad
OR_{\text{Hispanic}}=1.778,
$$

whereas the four-wheeled model estimates

$$
OR_{\text{Black}}=3.425,
\qquad
OR_{\text{Hispanic}}=3.258.
$$

Therefore, positive conditional associations for Black and Hispanic categories
remain in both datasets, but their estimated magnitude is substantially larger
among four-wheeled vehicle stops.

\begin{figure}[H]
\centering
\begin{minipage}{0.49\textwidth}
\centering
\includegraphics[width=\textwidth]{img/logistic_odds_ratios_2w.png}
\small (a) Two-wheeled vehicles
\end{minipage}
\hfill
\begin{minipage}{0.49\textwidth}
\centering
\includegraphics[width=\textwidth]{img/logistic_odds_ratios_4w.png}
\small (b) Four-wheeled vehicles
\end{minipage}
\caption{Estimated Logistic Regression odds ratios and 95\% confidence
intervals. Each estimate is conditional on the remaining predictors in the
corresponding model.}
\label{fig:logistic-odds-ratios}
\end{figure}
\FloatBarrier

Age provides one of the clearest examples of why the two vehicle populations
should not be treated as interchangeable. Among two-wheeled vehicles, the
Under-18 category has higher estimated odds than the 25--34 reference group
($OR=1.725$), while several middle-aged categories also remain above one.
Among four-wheeled vehicles, instead, 25--34 represents the highest-risk
reference category and the odds progressively decrease at older ages, reaching

$$
OR_{65+}=0.201.
$$

The conditional structure of arrest outcomes therefore differs materially
between the two vehicle contexts.

\subsection{Predictive Performance}

Despite the large sample sizes and numerous statistically significant
coefficients, predictive discrimination remains moderate.

\begin{table}[H]
\centering
\caption{Logistic Regression performance on the held-out test sets.}
\label{tab:logistic-performance}
\begin{tabular}{lrr}
\toprule
Metric & 2 Wheels & 4 Wheels \\
\midrule
ROC-AUC             & 0.6868 & 0.7352 \\
PR-AUC              & 0.1207 & 0.0774 \\
Log-Loss            & 0.2257 & 0.1373 \\
Brier Score         & 0.0576 & 0.0322 \\
Sensitivity         & 0.692  & 0.748  \\
Specificity         & 0.574  & 0.617  \\
Precision           & 0.098  & 0.065  \\
F1                  & 0.172  & 0.119  \\
MCC                 & 0.130  & 0.136  \\
Balanced Accuracy   & 0.633  & 0.683  \\
\bottomrule
\end{tabular}
\end{table}

The four-wheeled model discriminates arrests from non-arrests more effectively,
with an AUC of $0.7352$ compared with $0.6868$ for two-wheeled vehicles.
At the same time, precision remains low in both cases. This is expected from
the strong class imbalance and demonstrates an important distinction:
statistically clear conditional associations do not imply that individual arrest
outcomes are highly predictable.

% =============================================================================
% LASSO
% =============================================================================

\section{LASSO Logistic Regression}

\subsection{Why Introduce Regularization?}

The Logistic Regression contains many dummy coefficients generated by the
categorical predictors. With a very large sample, even comparatively weak
coefficients may become statistically significant. LASSO provides a useful
complement because it asks a different question:

\begin{quote}
Which of the patterns identified by the full logistic model remain important
when weak coefficients are progressively shrunk toward zero?
\end{quote}

For the binomial model, LASSO minimizes the negative log-likelihood together
with an $L_1$ penalty,

$$
-\ell(\boldsymbol{\beta})
+
\lambda\sum_{j=1}^{p}|\beta_j|.
$$

The penalty introduces bias but reduces model complexity and can set
coefficients exactly to zero. Consequently, LASSO is particularly useful here
as a robustness and parsimony check rather than merely as another classifier.

\subsection{Regularization and Model Parsimony}

The regularization parameter $\lambda$ was selected using cross-validated
ROC-AUC. Two standard solutions were investigated:
$\lambda_{\min}$, which maximizes cross-validated AUC, and
$\lambda_{1se}$, which selects a more strongly regularized model whose
performance remains within one standard error of the optimum.

\begin{table}[H]
\centering
\caption{LASSO cross-validation results.}
\label{tab:lasso-cv}
\begin{tabular}{lrrrr}
\toprule
Dataset &
$\lambda_{\min}$ &
CV-AUC$*{\min}$ &
$\lambda_{1se}$ &
CV-AUC$*{1se}$ \\
\midrule
2 Wheels & 0.000053 & 0.7394 & 0.000711 & 0.7373 \\
4 Wheels & 0.000048 & 0.7524 & 0.000256 & 0.7515 \\
\bottomrule
\end{tabular}
\end{table}

For two-wheeled vehicles, moving to the 1-SE solution reduces the number of
non-zero coefficients from $62$ to $53$, a reduction of approximately
$14.5\%$, while cross-validated AUC decreases by only $0.0021$.

The trade-off is even clearer for four-wheeled vehicles. The number of
non-zero coefficients falls from $63$ to $50$, approximately $20.6\%$,
while CV-AUC changes by only $0.0009$.

\begin{figure}[H]
\centering
\begin{minipage}{0.49\textwidth}
\centering
\includegraphics[width=\textwidth]{img/lasso_cv_2w.png}
\small (a) Two-wheeled vehicles
\end{minipage}
\hfill
\begin{minipage}{0.49\textwidth}
\centering
\includegraphics[width=\textwidth]{img/lasso_cv_4w.png}
\small (b) Four-wheeled vehicles
\end{minipage}
\caption{Cross-validated AUC as a function of the LASSO penalty. The
$\lambda_{\min}$ and $\lambda_{1se}$ solutions identify similar predictive
optima despite substantially different levels of regularization.}
\label{fig:lasso-cv}
\end{figure}
\FloatBarrier

The 1-SE solutions are therefore particularly useful for interpretation:
substantial simplification is achieved without materially reducing predictive
discrimination.

\subsection{Which Signals Survive Shrinkage?}

The most important result is not which individual dummy variable happens to
have the largest penalized coefficient, but whether the broad structure found
with Logistic Regression survives stronger regularization.

For two-wheeled vehicles, time of day remains dominant. At
$\lambda_{\min}$,

$$
OR_{\text{hour 7}}=0.043,
\qquad
OR_{\text{hour 6}}=0.097,
\qquad
OR_{\text{hour 8}}=0.137.
$$

At the more strongly regularized $\lambda_{1se}$, these estimates are shrunk
toward one but remain among the strongest coefficients:

$$
OR_{\text{hour 7}}=0.141,
\qquad
OR_{\text{hour 6}}=0.239,
\qquad
OR_{\text{hour 8}}=0.234.
$$

More importantly for the demographic analysis, the Black and Hispanic
coefficients are remarkably stable:

$$
OR_{\text{Black}}:
1.452 \rightarrow 1.473,
\qquad
OR_{\text{Hispanic}}:
1.776 \rightarrow 1.780.
$$

For four-wheeled vehicles, the same phenomenon is even more pronounced.
At $\lambda_{\min}$,

$$
OR_{\text{Black}}=3.377,
\qquad
OR_{\text{Hispanic}}=3.215,
$$

while at $\lambda_{1se}$,

$$
OR_{\text{Black}}=3.202,
\qquad
OR_{\text{Hispanic}}=3.057.
$$

The Female coefficient is similarly stable,

$$
OR_{\text{Female}}:
0.410 \rightarrow 0.420.
$$

The major age and hourly patterns also survive regularization. In particular,
the four-wheeled model retains strong negative coefficients for the morning
hours and progressively lower estimated odds among older age groups.

These results strengthen the interpretation of the Logistic Regression:
the principal temporal and demographic patterns are not eliminated when
coefficients with little incremental predictive contribution are aggressively
shrunk toward zero.

\begin{figure}[H]
\centering
\begin{minipage}{0.49\textwidth}
\centering
\includegraphics[width=\textwidth]{img/lasso_paths_2w.png}
\small (a) Two-wheeled vehicles
\end{minipage}
\hfill
\begin{minipage}{0.49\textwidth}
\centering
\includegraphics[width=\textwidth]{img/lasso_paths_4w.png}
\small (b) Four-wheeled vehicles
\end{minipage}
\caption{LASSO coefficient paths. Large temporal and demographic
coefficients remain separated from zero over the region between
$\lambda_{\min}$ and $\lambda_{1se}$, whereas weaker coefficients are
progressively removed.}
\label{fig:lasso-paths}
\end{figure}
\FloatBarrier

\subsection{Predictive Performance}

The test-set results confirm that the more parsimonious 1-SE solutions lose
very little predictive information.

\begin{table}[H]
\centering
\caption{Test performance for the preferred LASSO 1-SE models.}
\label{tab:lasso-performance}
\begin{tabular}{lrr}
\toprule
Metric & 2 Wheels & 4 Wheels \\
\midrule
ROC-AUC           & 0.6866 & 0.7340 \\
PR-AUC            & 0.1203 & 0.0763 \\
Log-Loss          & 0.2238 & 0.1374 \\
Brier Score       & 0.0576 & 0.0322 \\
Sensitivity       & 0.710  & 0.754  \\
Specificity       & 0.553  & 0.611  \\
Precision         & 0.097  & 0.064  \\
F1                & 0.170  & 0.118  \\
MCC               & 0.129  & 0.135  \\
Balanced Accuracy & 0.632  & 0.682  \\
\bottomrule
\end{tabular}
\end{table}

The results are almost indistinguishable from the corresponding
$\lambda_{\min}$ fits. The value of LASSO in this project therefore lies less
in improving discrimination than in demonstrating that the main findings are
stable under stronger regularization.

% =============================================================================
% GAM
% =============================================================================

\section{Generalized Additive Models}

\subsection{Why Move Beyond a Fully Categorical Logistic Model?}

The previous models confirm that hour and age contain substantial information
about arrest outcomes. However, the EDA also indicates that neither variable
follows a simple monotonic pattern.

Treating hour as 24 unrelated categories captures flexibility but ignores the
continuous and cyclic structure of the day. Similarly, age brackets impose
arbitrary boundaries on an intrinsically continuous variable.

A Generalized Additive Model addresses this limitation while retaining an
interpretable additive structure:

$$
\log\left(
\frac{P(\texttt{arrested}=1\mid X)}
     {1-P(\texttt{arrested}=1\mid X)}
\right)
=
\beta_0+\mathbf{X}\boldsymbol{\beta}
+f_1(\text{hour})
+f_2(\text{age}).
$$

Hour is represented using a cyclic cubic regression spline, reflecting the
continuity between 23:00 and 00:00, while age is represented by a cubic
regression spline.

The GAM therefore asks whether the temporal and age patterns already visible
in the EDA and Logistic Regression can be represented more naturally as smooth
non-linear relationships.

\subsection{Overall Fit}

For two-wheeled vehicles, the GAM obtains

$$
R^2_{\mathrm{adj}}=0.0391,
\qquad
\text{Deviance Explained}=9.61\%.
$$

For four-wheeled vehicles,

$$
R^2_{\mathrm{adj}}=0.0268,
\qquad
\text{Deviance Explained}=9.31\%.
$$

Thus, the included predictors explain a meaningful but still limited portion
of the total variability in the arrest outcome. This distinction is important:
large samples allow precise estimation of several associations, but the
available predictors still leave most individual-level variation unexplained.

\subsection{The Non-Linear Daily Cycle}

The hour smooth is highly significant in both datasets.

For two-wheeled vehicles,

$$
\text{edf}=7.52,
\qquad
\chi^2=1164,
\qquad
p<2\times10^{-16},
$$

while for four-wheeled vehicles,

$$
\text{edf}=7.925,
\qquad
\chi^2=6499,
\qquad
p<2\times10^{-16}.
$$

Effective degrees of freedom substantially larger than one confirm that the
relationship cannot be summarized adequately by a linear effect.

More importantly, the \textbf{shape is highly consistent across vehicle
types}. Conditional arrest propensity is comparatively high during late-night
and very early-morning hours, falls sharply toward its minimum around
07:00--08:00, and subsequently rises through the afternoon and evening.

For two-wheeled vehicles, the fitted smooth spans approximately $2.3$
log-odds units between its minimum and maximum, corresponding to roughly a
ten-fold difference in estimated odds across the daily curve, conditional on
the remaining predictors. For four-wheeled vehicles the corresponding range is
approximately $1.9$ log-odds units, or around $6.7$ times in the odds.

\begin{figure}[H]
\centering
\begin{minipage}{0.49\textwidth}
\centering
\includegraphics[width=\textwidth]{img/gam_hour_2w.png}
\small (a) Two-wheeled vehicles
\end{minipage}
\hfill
\begin{minipage}{0.49\textwidth}
\centering
\includegraphics[width=\textwidth]{img/gam_hour_4w.png}
\small (b) Four-wheeled vehicles
\end{minipage}
\caption{GAM partial smooth for hour of day. Both datasets exhibit a strong
non-linear daily cycle, with minimum conditional arrest propensity during
the morning and substantially higher values at night.}
\label{fig:gam-hour}
\end{figure}
\FloatBarrier

This result is particularly useful for the overall narrative because the same
temporal structure is visible at three different analytical stages:
descriptively in the EDA, conditionally through Logistic Regression, and
continuously through the GAM.

\subsection{Age Reveals a Stronger Difference Between Vehicle Types}

The age smooth is also highly significant in both models.

For two-wheeled vehicles,

$$
\text{edf}=4.66,
\qquad
\chi^2=108,
\qquad
p<2\times10^{-16},
$$

whereas for four-wheeled vehicles,

$$
\text{edf}=4.878,
\qquad
\chi^2=4057,
\qquad
p<2\times10^{-16}.
$$

Unlike hour, however, the shape is not stable across datasets.

For two-wheeled vehicles, conditional arrest propensity is relatively elevated
among the youngest observations, reaches a local minimum around age 28--30,
remains comparatively flat through much of middle age, rises toward a local
maximum around age 60--62, and subsequently declines among the oldest
observations.

The four-wheeled pattern is substantially different. The smooth reaches its
maximum around age 27--30 and then declines progressively. It crosses the
zero-centered smooth around age 38--40 and becomes increasingly negative,
with a particularly strong decline after approximately age 50--55.

\begin{figure}[H]
\centering
\begin{minipage}{0.49\textwidth}
\centering
\includegraphics[width=\textwidth]{img/gam_age_2w.png}
\small (a) Two-wheeled vehicles
\end{minipage}
\hfill
\begin{minipage}{0.49\textwidth}
\centering
\includegraphics[width=\textwidth]{img/gam_age_4w.png}
\small (b) Four-wheeled vehicles
\end{minipage}
\caption{GAM partial smooth for motorist age. In contrast to the consistent
hourly cycle, the age relationship differs substantially between the two
vehicle populations.}
\label{fig:gam-age}
\end{figure}
\FloatBarrier

This reinforces the decision to analyze the two vehicle populations
separately. A single pooled age effect would conceal substantially different
conditional structures.

\subsection{Demographic Associations Remain After the Flexible Specification}

Introducing smooth effects for age and hour does not eliminate the principal
demographic associations.

In the two-wheeled GAM,

$$
OR_{\text{Black}}=1.563,
\qquad
OR_{\text{Hispanic}}=1.900,
\qquad
OR_{\text{Female}}=0.309.
$$

In the four-wheeled GAM,

$$
OR_{\text{Black}}=3.349,
\qquad
OR_{\text{Hispanic}}=3.185,
\qquad
OR_{\text{Female}}=0.403.
$$

The direction and approximate magnitude of these effects closely resemble
those found by Logistic Regression and LASSO. The demographic disparities are
therefore not removed simply by modelling age and time of day more flexibly.

\subsection{Predictive Performance}

\begin{table}[H]
\centering
\caption{GAM performance on the held-out test sets.}
\label{tab:gam-performance}
\begin{tabular}{lrr}
\toprule
Metric & 2 Wheels & 4 Wheels \\
\midrule
ROC-AUC           & 0.6897 & 0.7347 \\
PR-AUC            & 0.1247 & 0.0790 \\
Log-Loss          & 0.2235 & 0.1400 \\
Brier Score       & 0.0570 & 0.0330 \\
Sensitivity       & 0.697  & 0.759  \\
Specificity       & 0.576  & 0.607  \\
Precision         & 0.099  & 0.066  \\
F1                & 0.173  & 0.121  \\
MCC               & 0.133  & 0.137  \\
Balanced Accuracy & 0.637  & 0.683  \\
\bottomrule
\end{tabular}
\end{table}

The additional flexibility produces only limited predictive improvement over
Logistic Regression. Its main contribution is therefore not a dramatic
increase in AUC, but a more informative representation of the non-linear
temporal and age-related structure.

% =============================================================================
% NAIVE BAYES
% =============================================================================

\section{Naive Bayes}

\subsection{Why a Generative Model?}

Logistic Regression, LASSO, and GAM directly model the conditional probability

$$
P(Y\mid X).
$$

Naive Bayes approaches the classification problem from the opposite direction:
it models the distribution of the predictors within each outcome class,
$P(X\mid Y)$, and combines these distributions with the class priors through
Bayes' theorem.

For class $k$,

$$
P(Y=k\mid X)
=
\frac{
P(Y=k)\prod_{j=1}^{p}P(X_j\mid Y=k)
}{
\sum_h P(Y=h)\prod_{j=1}^{p}P(X_j\mid Y=h)
}.
$$

The price of this simplification is the strong assumption that predictors are
conditionally independent within each class. However, this makes Naive Bayes a
useful robustness check: it asks whether the distinction between arrests and
non-arrests is also visible from a fundamentally different modelling
perspective.

\subsection{Class Priors and Conditional Distributions}

The estimated prior probability of arrest differs substantially between the
two datasets:

$$
P(Y=1)=0.0538
$$

for two-wheeled vehicles, compared with

$$
P(Y=1)=0.0325
$$

for four-wheeled vehicles.

Thus, before considering any predictor information, both models strongly favor
the non-arrested class.

The class-conditional distributions nevertheless contain useful structure.

For two-wheeled vehicles, Hispanic motorists account for approximately
$63.4\%$ of arrests but $46.9\%$ of non-arrests, while Black motorists
account for $25.6\%$ and $20.4\%$, respectively. The Bronx accounts for
approximately $34.9\%$ of arrests compared with only $14.6\%$ of
non-arrests. Late-evening and nighttime observations are also substantially
more represented among arrests than morning observations.

For four-wheeled vehicles, the ethnicity contrast is even clearer. Black
motorists constitute approximately $47.0\%$ of arrested observations but
$30.0\%$ of non-arrests, while Hispanic motorists account for $37.5\%$
and $26.0\%$, respectively. White motorists show the opposite pattern:
approximately $6.9\%$ of arrests but $20.3\%$ of non-arrests.

Age also separates the four-wheeled classes more clearly. The 25--34 category
represents approximately $40.9\%$ of arrests but only $29.0\%$ of
non-arrests, whereas the older age groups are progressively less represented
among arrests.

These quantities must not be confused with arrest probabilities. For example,

$$
P(\text{Black}\mid Y=1)
$$

describes the composition of the arrested class and is not the same quantity as

$$
P(Y=1\mid\text{Black}).
$$

Naive Bayes combines all such class-specific evidence to construct the final
posterior probability.

\begin{figure}[H]
\centering
\begin{minipage}{0.49\textwidth}
\centering
\includegraphics[width=\textwidth]{img/nb_conditional_2w.png}
\small (a) Two-wheeled vehicles
\end{minipage}
\hfill
\begin{minipage}{0.49\textwidth}
\centering
\includegraphics[width=\textwidth]{img/nb_conditional_4w.png}
\small (b) Four-wheeled vehicles
\end{minipage}
\caption{Selected class-conditional probabilities from Naive Bayes.
Ethnicity, hour, and geographical context exhibit visible differences
between the arrested and non-arrested classes.}
\label{fig:nb-conditional}
\end{figure}
\FloatBarrier

\subsection{Predictive Performance}

\begin{table}[H]
\centering
\caption{Naive Bayes performance on the held-out test sets.}
\label{tab:nb-performance}
\begin{tabular}{lrr}
\toprule
Metric & 2 Wheels & 4 Wheels \\
\midrule
ROC-AUC           & 0.6818 & 0.7322 \\
PR-AUC            & 0.1169 & 0.0761 \\
Log-Loss          & 0.2306 & 0.1394 \\
Brier Score       & 0.0584 & 0.0325 \\
Sensitivity       & 0.748  & 0.762  \\
Specificity       & 0.510  & 0.598  \\
Precision         & 0.093  & 0.063  \\
F1                & 0.166  & 0.116  \\
MCC               & 0.126  & 0.133  \\
Balanced Accuracy & 0.629  & 0.680  \\
\bottomrule
\end{tabular}
\end{table}

Naive Bayes performs slightly worse than the discriminative models but remains
within the same general performance range. The result is informative:
meaningful arrest-related structure can be recovered under a very different
modelling framework, but class separation remains incomplete.

% =============================================================================
% MODEL COMPARISON
% =============================================================================

\section{Overall Model Comparison}

The most useful comparison is not whether one model improves AUC by a few
thousandths, but whether the same substantive patterns remain visible under
different modelling assumptions.

\begin{table}[H]
\centering
\caption{Summary of held-out predictive performance across modelling
approaches. LASSO results refer to the preferred 1-SE specification.}
\label{tab:model-comparison}
\begin{tabular}{llrrrr}
\toprule
Dataset & Model & ROC-AUC & PR-AUC & MCC & Bal. Acc. \\
\midrule
2 Wheels & Logistic    & 0.6868 & 0.1207 & 0.130 & 0.633 \\
& LASSO       & 0.6866 & 0.1203 & 0.129 & 0.632 \\
& GAM         & 0.6897 & 0.1247 & 0.133 & 0.637 \\
& Naive Bayes & 0.6818 & 0.1169 & 0.126 & 0.629 \\
\midrule
4 Wheels & Logistic    & 0.7352 & 0.0774 & 0.136 & 0.683 \\
& LASSO       & 0.7340 & 0.0763 & 0.135 & 0.682 \\
& GAM         & 0.7347 & 0.0790 & 0.137 & 0.683 \\
& Naive Bayes & 0.7322 & 0.0761 & 0.133 & 0.680 \\
\bottomrule
\end{tabular}
\end{table}

Three conclusions emerge.

First, \textbf{no modelling approach dramatically dominates the others}.
Logistic Regression, LASSO, and GAM produce almost identical discrimination,
while Naive Bayes performs only slightly worse. Increasing model complexity
therefore does not radically change the amount of predictive information
available.

Second, \textbf{four-wheeled arrest outcomes are consistently easier to rank}
than two-wheeled outcomes. ROC-AUC is approximately $0.73$--$0.74$ for
all four-wheeled models and approximately $0.68$--$0.69$ for two-wheeled
models.

Third, the apparently larger two-wheeled PR-AUC should not be interpreted as
contradicting the previous result. Arrest prevalence is approximately $5\%-
6\%$ in the two-wheeled population but only $3\%-4\%$ in the
four-wheeled population, and the no-skill PR baseline changes accordingly.
Relative to their respective baselines, both populations contain useful but
limited positive-class information.

Across all models, precision remains low. Consequently, the models are more
useful for \textbf{conditional interpretation and relative risk ranking} than
for high-confidence individual arrest classification.

% =============================================================================
% RESEARCH QUESTIONS
% =============================================================================

\chapter{Research Questions and Findings}

The preceding analyses can now be combined to answer the research questions
introduced by the project. The objective of this chapter is not to introduce
additional models, but to synthesize evidence from the EDA and the different
statistical-learning approaches.

% -----------------------------------------------------------------------------
\section{RQ1: Are Arrest Outcomes Distributed Uniformly Across Recorded Stops?}

\textbf{Answer: No.} The descriptive analysis provides clear evidence that
arrest outcomes are heterogeneous across demographic, temporal, and
geographical characteristics.

Among two-wheeled stops, the observed arrest rate is $12.1\%$ in the Bronx
but only approximately $3.2\%$ in Manhattan. For four-wheeled stops the
geographical range is narrower but remains substantial, from approximately
$4.8\%$ in the Bronx to $2.5\%$ in Queens.

Time of day produces an especially consistent pattern. At approximately
07:00, arrest rates fall to $3.2\%$ for two-wheeled vehicles and $0.6\%$
for four-wheeled vehicles. By 03:00, they rise to approximately $10.9\%$
and $5.3\%$, respectively.

Demographic differences are also substantial. For two-wheeled vehicles,
Hispanic and Black motorists have observed arrest rates of $7.24\%$ and
$6.86\%$, compared with $3.23\%$ for White motorists. For four-wheeled
vehicles, the corresponding rates are $4.66\%$, $5.02\%$, and $1.15\%$.

The EDA therefore rejects the hypothesis that arrest outcomes are uniformly
distributed across the observed characteristics. These are, however,
\textbf{unadjusted differences}; they do not establish which characteristics
retain an association after accounting for the others.

% -----------------------------------------------------------------------------
\section{RQ2: Do Demographic Disparities Persist After Adjustment?}

\textbf{Answer: Yes.} Several of the demographic differences identified by the
EDA remain after conditioning on temporal, geographical, and other demographic
characteristics.

The clearest evidence comes from reported ethnicity. In the Logistic
Regression for two-wheeled vehicles,

$$
OR_{\text{Black}}=1.451,
\qquad
OR_{\text{Hispanic}}=1.778,
$$

relative to the White reference category.

For four-wheeled vehicles the estimated associations are much larger:

$$
OR_{\text{Black}}=3.425,
\qquad
OR_{\text{Hispanic}}=3.258.
$$

These findings are not dependent on a single model. In the four-wheeled LASSO
1-SE model, the corresponding penalized odds ratios remain approximately

$$
3.202
\quad\text{and}\quad
3.057,
$$

while the GAM estimates

$$
3.349
\quad\text{and}\quad
3.185.
$$

The same consistency is visible for sex. Female observations have conditional
odds ratios of $0.300$ and $0.408$ in the two Logistic models, and the
direction and magnitude remain similar under LASSO and GAM.

Therefore, demographic disparities do not disappear after controlling for the
observed context included in the models. In particular, the Black and Hispanic
coefficients remain large and resistant to regularization in the four-wheeled
analysis.

These results demonstrate \textbf{persistent conditional statistical
disparities}. They do not, by themselves, identify the mechanism responsible
for those disparities.

% -----------------------------------------------------------------------------
\section{RQ3: Are the Relationships with Hour and Age Non-Linear?}

\textbf{Answer: Yes.} The GAM provides strong statistical and graphical
evidence that both relationships are non-linear.

For hour of day,

$$
\text{edf}_{2W}=7.52,
\qquad
\text{edf}_{4W}=7.925,
$$

with $p<2\times10^{-16}$ in both models. The smooth functions reproduce the
same broad cycle observed during the EDA: elevated arrest propensity during
late-night and early-morning hours, a pronounced minimum around 07:00--08:00,
and a progressive increase during the afternoon and evening.

The consistency of this curve across vehicle types makes time of day one of
the most stable findings in the entire analysis.

Age is also non-linear,

$$
\text{edf}_{2W}=4.66,
\qquad
\text{edf}_{4W}=4.878,
$$

again with $p<2\times10^{-16}$.

Unlike hour, however, the age curves differ substantially between vehicle
types. The four-wheeled smooth reaches its highest value around ages 27--30
and subsequently declines, particularly after age 50. The two-wheeled curve
has a more irregular structure, with elevated values among the youngest
observations and another local increase at later ages.

The hypothesis that age and hour can be represented adequately by simple
linear effects is therefore rejected.

% -----------------------------------------------------------------------------
\section{RQ4: Do Two-Wheeled and Four-Wheeled Stops Exhibit the Same Patterns?}

\textbf{Answer: Only partially.}

Some results are remarkably consistent. Both datasets exhibit:

\begin{itemize}
\item a strong late-night versus morning arrest pattern;
\item lower conditional odds for female motorists;
\item positive conditional associations for Black and Hispanic categories
relative to White;
\item substantial geographical heterogeneity;
\item only moderate individual-level predictability.
\end{itemize}

However, the magnitude and sometimes even the direction of specific
associations differ considerably.

The ethnicity contrast is substantially stronger among four-wheeled vehicles.
For example,

$$
OR_{\text{Black}}:
1.451\;(2W)
\quad\text{vs.}\quad
3.425\;(4W),
$$

and

$$
OR_{\text{Hispanic}}:
1.778\;(2W)
\quad\text{vs.}\quad
3.258\;(4W).
$$

Age provides an even clearer difference. In the Logistic Regression, Under-18
two-wheeled observations have higher estimated odds than the 25--34 reference
group,

$$
OR=1.725,
$$

whereas the corresponding four-wheeled estimate is

$$
OR=0.756.
$$

Similarly, the four-wheeled model exhibits a strong progressive reduction at
older ages, reaching $OR_{65+}=0.201$, while the corresponding two-wheeled
65+ coefficient is not statistically distinguishable from the 25--34
reference group.

The GAM smooths confirm that this is not merely an artifact of the age
brackets.

Consequently, the data support the decision to maintain separate analyses.
There appears to be a common temporal component, but the demographic
structure of arrest outcomes is strongly dependent on the vehicle context.

% -----------------------------------------------------------------------------
\section{RQ5: How Predictable Is an Arrest from the Available Information?}

\textbf{Answer: Arrest is moderately predictable as a ranking problem, but not
with high reliability at the individual classification level. Predictability
is consistently greater for four-wheeled stops.}

Across two-wheeled models, ROC-AUC ranges approximately from

$$
0.682 \text{ to } 0.690,
$$

whereas for four-wheeled models it remains around

$$
0.732 \text{ to } 0.735.
$$

Thus, all models perform better than random ranking, but none achieves strong
class separation.

The PR-AUC results must be interpreted relative to the different class
prevalences. For two-wheeled vehicles, PR-AUC is approximately
$0.117$--$0.125$ against a no-skill baseline of approximately $0.063$.
For four-wheeled vehicles it is approximately $0.076$--$0.079$ against a
much lower baseline of approximately $0.034$.

The principal limitation is precision. Around the selected operating points,
precision is approximately $9\%-10\%$ for two-wheeled models and
$6\%-7\%$ for four-wheeled models. Consequently, most observations
classified as arrests are false positives even though sensitivity is generally
around $70\%-76\%$.

The models therefore contain genuine information about relative arrest
propensity, but the observed characteristics are not sufficient to determine
individual arrest outcomes with high certainty.

The GAM fit reinforces this conclusion from another perspective: only about
$9.61\%$ and $9.31\%$ of deviance is explained in the two vehicle
populations. A substantial part of the outcome therefore depends on
information that is either unobserved or not represented by the available
predictors.

% -----------------------------------------------------------------------------
\section{RQ6: Do the Results Demonstrate Systematic or Excessive Enforcement?}

\textbf{Answer: The results provide evidence of persistent systematic
statistical disparities in post-stop arrest outcomes, but they are not
sufficient to establish systematic discriminatory or excessive enforcement.}

There is substantial empirical evidence that the observed differences are not
purely descriptive artifacts. For example, among four-wheeled vehicle stops,
the raw arrest rate is $5.02\%$ for Black motorists and $4.66\%$ for
Hispanic motorists compared with $1.15\%$ for White motorists. After
adjustment, the Logistic Regression still estimates odds ratios of $3.425$
and $3.258$, respectively. LASSO and GAM produce very similar conditional
relationships.

Similarly, significant geographical, sex-related, age-related, and temporal
differences remain after adjustment, and several of these patterns are
recovered under multiple modelling assumptions.

This constitutes evidence that the disparities are
\textbf{persistent and structured within the recorded stop data}.

However, three different statements must be distinguished:

$$
\boxed{
\text{Observed disparity}
\neq
\text{causal demographic effect}
\neq
\text{proof of discriminatory enforcement}
}
$$

The dataset contains only vehicle stops that have already occurred. It
therefore allows the analysis of

$$
P(\text{arrest}\mid\text{recorded stop},X),
$$

but not the probability that a motorist is stopped in the first place.

Furthermore, the records do not contain complete information about the reason
for the stop, the legal circumstances leading to the arrest, the behavior
during the encounter, or other operational factors that may influence the
decision to arrest. Consequently, important confounding information is not
available to the models.

The strongest conclusion supported by the available evidence is therefore the
following:

\begin{quote}
Among recorded NYPD vehicle stops, arrest outcomes exhibit persistent
demographic, temporal, and geographical disparities. Several of these
differences remain large after adjustment for the observed characteristics and
are reproduced across alternative statistical models, particularly in the
four-wheeled population. However, the available data are insufficient to
determine whether these conditional disparities are caused by discriminatory
or excessive police enforcement.
\end{quote}

This distinction is central to the interpretation of the project. The analysis
does not confirm a predetermined claim of police bias, nor does it provide
evidence that the observed disparities are innocuous. Instead, it identifies
which disparities are present, quantifies their magnitude, tests their
robustness across modelling assumptions, and establishes the boundary beyond
which the available data cannot support stronger causal conclusions.

% =============================================================================
% FINAL SYNTHESIS
% =============================================================================

\section{Synthesis}

Taken together, the exploratory and modelling analyses produce a coherent
picture.

The most stable finding is temporal: arrest propensity follows a pronounced
daily cycle in both vehicle populations, with substantially lower conditional
odds during morning hours and higher values during the night. The agreement
between raw arrest rates, Logistic Regression coefficients, LASSO retention,
Naive Bayes conditional distributions, and GAM smooths makes this the most
robust cross-dataset pattern.

Demographic variables also contain persistent information, but their effects
are considerably more context-dependent. Reported ethnicity and sex remain
associated with arrest after adjustment and regularization, while the magnitude
of the ethnicity coefficients is substantially greater among four-wheeled
vehicle stops. Age exhibits an even stronger vehicle-specific structure.

At the same time, predictive performance places an important limit on the
interpretation. The available predictors produce moderate ranking ability but
low individual precision, and more flexible modelling does not substantially
raise this performance ceiling.

The principal contribution of the modelling stage is therefore not the
identification of a single ``best'' classifier. It is the convergence of
several modelling approaches toward the same broader conclusion:
\textbf{post-stop arrest outcomes exhibit structured and persistent
disparities, but much of the individual outcome remains unexplained by the
information available in the dataset.}
